\section{\rev{결과}}
\label{sec:results}

\begin{revision}

결과는 두 단계로 제공한다. 첫째, 작성ㆍ기호 입력에서 이전 의무의 기억이 필요한지, 그것만으로 충분한지, 부족하다면 무엇을 바꿔야 하는지를 확인하고(\ref{sec:regression-results}--\ref{sec:agreement-results}절), 같은 질문을 장면 경계로 넓힌 뒤(\ref{sec:boundary-results}절), 이 결과가 사례 선별에 의존하지 않는지 전수 대조로 확인한다(\ref{sec:bounded-results}절). 둘째, 실제 CoC에서 이 검사를 평가하기 어려운 이유와 기록 궤적과의 대응을 보인다(\ref{sec:human-results}--\ref{sec:trajectory-results}절).

\subsection{기존 40쌍의 규칙 실행 결과}
\label{sec:regression-results}
먼저 이전 의무를 기억하는 것이 필요한지를 문제 유형마다 한 항목만 바꾼 회귀 쌍으로 확인했다. 표~\ref{tab:regression}에서 두 검사의 차이는 이전 사건의 정보가 있어야 판단할 수 있는 P1과 P4에서만 나타났고, 현재 사건에 판정 근거가 있는 P2와 P3은 두 검사가 모두 판정했다. 기준 사례에 확정 모순을 표시한 경우는 없었다. 합계 40/40과 20/40을 일반 오류 검출 정확도로 해석하지 않는다. 기존 검사기의 P2 표시는 진행 사건의 해제 조건 관측값 \texttt{release=false}에서, P4 표시는 같은 대상의 정지ㆍ양보 반복과 \texttt{release=true}가 함께 주어진 처리 변형에서 발생하므로, 명시적 해제 명령을 검사하는 $C_2$나 사람의 실제 조기 해제를 확인한 결과와는 다르다.

\begin{table}[t]
\centering
\color{revisionblue}
\bitablecaption{기존 규칙 실행 회귀시험: 주석 문제와 처리 오류를 분리.}{Legacy rule-execution regression, separating annotation-related rules and processing faults.}
\label{tab:regression}
\footnotesize
\begin{tabularx}{\columnwidth}{Ycc}
\toprule
시험 & 전체 이력 & 현재 사건 \\
\midrule
P1 미해제 의무와 진행 & 10/10 & 0/10 \\
P2 조기 해제 & 10/10 & 10/10 \\
P3 선행조건 미충족 진행 & 10/10 & 10/10 \\
P4 해제 후 잔존 & 10/10 & 0/10 \\
기준 사례 확정 모순 표시 & 0/40 & 0/40 \\
\bottomrule
\end{tabularx}
\end{table}

\subsection{최초 144개 작성 시나리오}
\label{sec:first-results}
기억이 필요하다는 것만으로 충분하지는 않았다. 의무가 여럿이거나 사건 간격이 길거나 해제 근거가 빠진 144개 시나리오에서는 전체 이력을 기억하는 기존 검사기도 96/144만 기대 판정과 일치했다(표~\ref{tab:first-results}). 이 집계는 기존 경로의 최종 출력 \texttt{CONTRADICTION}, \texttt{UNKNOWN}, \texttt{CONSISTENT}를 각각 확정 모순, 근거 부족, 양립에 대응시켜 센 것이다. 이하 \texttt{UNKNOWN}은 ``미상 출력''이라 부르며, 구문 모호성ㆍ복수 의무ㆍ대기 중 미상까지 누적하므로 4장의 해제 근거 부족과 같은 의미는 아니다.

\begin{table}[t]
\centering
\color{revisionblue}
\bitablecaption{고정된 144개 작성 시나리오의 최초 결과. 간격별 분모는 각각 48이며, $^\dagger$는 작성 프레임을 단일 의무로 투영한 입력이다. 조건 연결 모델의 재시험은 포함하지 않는다.}{First-run results on all 144 authored cases, with 48 cases per gap; $^\dagger$ marks single-obligation projected input. The condition-linked model retest is excluded.}
\label{tab:first-results}
\scriptsize
\setlength{\tabcolsep}{2.5pt}
\begin{tabularx}{\columnwidth}{Yccccc}
\toprule
경로 & 세 분류 & 최초 위반 & 간격 0 & 간격 2 & 간격 6 \\
\midrule
기존 전체 이력$^\dagger$ & 96/144 & 54/72 & 36/48 & 30/48 & 30/48 \\
현재 사건$^\dagger$ & 72/144 & 36/72 & 24/48 & 24/48 & 24/48 \\
이전 1개$^\dagger$ & 84/144 & 42/72 & 36/48 & 24/48 & 24/48 \\
이전 3개$^\dagger$ & 90/144 & 48/72 & 36/48 & 30/48 & 24/48 \\
이전 5개$^\dagger$ & 90/144 & 48/72 & 36/48 & 30/48 & 24/48 \\
조건별 FSM & 144/144 & 72/72 & 48/48 & 48/48 & 48/48 \\
기존 자동 텍스트 & 36/144 & 0/72 & 12/48 & 12/48 & 12/48 \\
기존 UPPAAL$^\dagger$ & 96/144 & 54/72 & 36/48 & 30/48 & 30/48 \\
\bottomrule
\end{tabularx}
\end{table}


기존 경로의 실패는 잘못된 확정 모순이 아니라 판정하지 못한 미상 출력이었다. 기대 양립 사례에서 모순 오탐은 없었지만 대부분을 미상 출력으로 남겼으므로, 오탐이 없다는 사실만으로 성능을 판단할 수 없다. 사건 창은 간격이 길어질수록 이전 증거를 잃었지만, 창을 늘려도 전체 이력의 결과를 넘지 못했다. 즉 더 긴 문맥만으로는 복수 의무와 미상 누적의 문제가 해결되지 않는다. FSM은 144/144, 최초 위반 72/72였다. UPPAAL이 단순 비교기보다 정확하다는 근거는 없다. 원문만 입력한 기존 자동 텍스트 경로는 모든 사례에서 미상 출력을 반환했으므로 구조화 입력의 성공을 자동 추출의 성공으로 볼 수 없다.

\subsection{실패 분석과 개발 후 동일 사례 재시험}
\label{sec:retest-results}
\ref{sec:first-results}절의 불일치 48건은 기억의 길이가 아니라 의무를 표현하는 방식과 미상을 다루는 정책에서 생겼으며 세 유형으로 나뉜다(표~\ref{tab:failures}). 동결된 기존 검사기는 활성 의무 하나와 6개 상태(\texttt{INACTIVE}, \texttt{ACTIVE}, \texttt{SATISFIED\_WAIT\_RELEASE}, \texttt{RELEASED}, \texttt{VIOLATED}, \texttt{UNKNOWN})로 동작한다. 복수 의무는 하나의 채널에 담을 수 없어 미상으로 처리되었고, 대기 중의 해제 미상은 누적되어 뒤에 명시적으로 해제되어도 지워지지 않았다. 이는 작성한 의미와 기존 표현ㆍ정책의 불일치이며 자연 오류가 아니다.

조건 연결 모델은 의무별 활성 상태와 조건별 증거값을 나누어 복수 의무를 각각 판정한다. 또한 ``해제 정보를 기다리는 단계''와 ``해제 근거 없이 진행한 단계''를 구분해 후자에서만 근거 부족을 기록하고, 한 번 기록한 문제는 이후의 해제로 지우지 않는다(목록~\ref{lst:condition-step}). 같은 144건의 재시험에서 이 모델은 판정뿐 아니라 최초 위반과 문제 위치ㆍ관련 의무ㆍ사유를 모두 기대대로 돌려주었다(표~\ref{tab:retest}). 다만 같은 실패 사례를 모델 개발에 사용했으므로 개발에 쓰지 않은 독립 시험의 결과가 아니며, 자동 추출의 개선을 뜻하지도 않는다.

\begin{table}[t]
\centering
\color{revisionblue}
\bitablecaption{48개 최초 불일치와 보완의 대상.}{All 48 initial mismatches and the corresponding refinement.}
\label{tab:failures}
\scriptsize
\begin{tabularx}{\columnwidth}{L{.36\columnwidth}cY}
\toprule
유형 & 건수 & 보완 내용 \\
\midrule
복수 일부 해제 & 18 & C의 조건과 의무를 A와 별도로 유지 \\
복수 모두 해제 & 18 & 각 조건의 참 증거로 각각 해제 \\
중간 대기 후 정상 해제 & 12 & 미상 대기와 근거 없이 진행한 사건 구분 \\
\bottomrule
\end{tabularx}
\end{table}

\begin{table}[t]
\centering
\color{revisionblue}
\bitablecaption{개발 후 동일 사례 재시험: 조건 연결 모델(독립 시험 아님).}{Post-development same-case retest of the condition-linked UPPAAL model, not an independent test.}
\label{tab:retest}
\footnotesize
\begin{tabularx}{\columnwidth}{Yc}
\toprule
대조 항목 & 기대 기록과 일치 \\
\midrule
최종 세 분류 & 144/144 \\
확정 모순의 최초 위반 사건 & 72/72 \\
문제 위치(확정 모순+근거 부족) & 108/108 \\
관련 의무 & 108/108 \\
문제 사유 & 108/108 \\
\bottomrule
\end{tabularx}
\end{table}

\subsection{기존 Python--UPPAAL 출력 대조}
\label{sec:agreement-results}
\ref{sec:first-results}절의 불일치를 해석하기 위해 기존 두 구현의 출력을 대조했다. 같은 규칙을 구현한 기존 Python 검사기와 UPPAAL 모델은 모든 사례에서 판정, 오류 종류, 최초 위치가 같았다(표~\ref{tab:implementation}). 미상 사유가 일부 달랐던 것은 최초 반례의 접두 경로에 뒤에서 누적된 사유가 담기지 않기 때문이었고, 종료 상태의 추가 질의로 모두 일치했다. 따라서 이 차이는 출력 범위의 문제이며 \ref{sec:retest-results}절의 의미 불일치와는 별개다.

\begin{table}[t]
\centering
\color{revisionblue}
\bitablecaption{기존 Python--UPPAAL 구현 일치성 및 종료 질의 진단.}{Legacy Python--UPPAAL output agreement and terminal-query diagnosis.}
\label{tab:implementation}
\footnotesize
\begin{tabularx}{\columnwidth}{Yc}
\toprule
대조 항목 & 일치 사례 \\
\midrule
최종 판정 / 오류 종류 / 최초 위치 & 각 144/144 \\
원 반례 접두 경로의 미상 사유 & 120/144 \\
종료 상태 추가 질의의 미상 사유 & 144/144 \\
\bottomrule
\end{tabularx}
\end{table}

\subsection{장면 경계 결과(81건)}
\label{sec:boundary-results}
같은 질문은 장면 경계에서 다시 생긴다. 경계에서 의무와 증거를 버리면 장면을 넘어 이어지는 의무 위반을 판정할 수 있는가? 표~\ref{tab:boundary-results}에서 경계에서 상태를 유지하는 경계 확장 모델(이하 유지 모델)은 모든 사례를 기대대로 판정했지만, 경계마다 의무를 버리는 초기화 모델은 확정 모순과 근거 부족 사례의 대부분을 양립으로 판정했다. 예를 들어 세 장면에 걸쳐 마지막 진행 시점에 보행자 조건만 거짓이면 유지 모델은 남은 보행자 의무로 확정 모순을 기록하지만 초기화 모델은 양립으로 끝난다. 확정 모순과 근거 부족이 함께 생긴 사례에서는 두 기록을 보존하고 확정 모순을 최종 판정으로 택했다. 근거 만료 뒤의 진행은 근거 부족이 되고, 재활성화된 의무에는 과거 해제값을 쓰지 않는다. 재활성화 사례는 초기화 모델도 판정하므로 상태 유지의 이점으로 세지 않았다.

\begin{table}[t]
\centering
\color{revisionblue}
\bitablecaption{81건 경계 시험의 기대 판정별 출력. 초기화는 기억 제거 비교다.}{Outputs by expected class in the 81-case boundary study. Reset is a memory ablation.}
\label{tab:boundary-results}
\scriptsize
\begin{tabularx}{\columnwidth}{Ycccc}
\toprule
기대 판정 & 유지 일치 & 초기화 확정 모순 & 초기화 근거 부족 & 초기화 양립 \\
\midrule
확정 모순 44 & 44 & 3 & 0 & 41 \\
근거 부족 29 & 29 & 0 & 3 & 26 \\
양립 8 & 8 & 0 & 0 & 8 \\
\bottomrule
\end{tabularx}
\end{table}

\subsection{유한 범위 전수 적합성 결과}
\label{sec:bounded-results}
\ref{sec:retest-results}절과 \ref{sec:boundary-results}절의 결과는 모두 설계자가 만든 사례에서 나왔다. 사례 선별의 영향을 배제하기 위해 선언한 범위의 모든 사건열에서 구현과 수식 판정을 비교한 결과, 두 모델 모두 불일치가 없었다(표~\ref{tab:bounded-results}). 같은 비교에 P4 해제 고장을 넣은 대조군에서는 상태 불일치가 다수 나타나, 이 대조가 처리 오류를 구분할 수 있음을 확인했다. 표의 사건열 수는 열거한 조합의 수이며 자연 표본 수나 탐지 정확도의 분모가 아니다.

\begin{table}[t]
\centering
\color{revisionblue}
\bitablecaption{고정한 갱신 코드와 별도 수식 판정기의 유한 범위 전수 대조.}{Bounded exhaustive comparison of frozen step code and a separate formula oracle.}
\label{tab:bounded-results}
\scriptsize
\begin{tabularx}{\columnwidth}{Yrr}
\toprule
모델ㆍ의무 수ㆍ길이 상한 & 사건열 & 불일치 \\
\midrule
조건 연결ㆍ1개ㆍ3 & 378,504 & 0 \\
조건 연결ㆍ2개ㆍ3 & 34,117,524 & 0 \\
경계 확장ㆍ1개ㆍ3 & 348,504 & 0 \\
경계 확장ㆍ2개ㆍ2 & 198,396 & 0 \\
\bottomrule
\end{tabularx}
\end{table}

이 시험은 기대 결과에 맞는 사례만 고르는 절차를 없애고 두 계산 경로를 비교한다. 그러나 수식과 구현이 공유하는 명세의 오류는 함께 남을 수 있다. 유한 상한 밖의 사건열이나 자연 문장의 의미 정답은 입증하지 않는다.

\subsection{자연 검토의 일치도}
\label{sec:human-results}
작성 시나리오와 달리 실제 자동 생성 CoC에는 독립적인 기대 판정을 확보하지 못했다. 네 검토자의 원 판정도 크게 엇갈렸다.

\begin{table}[t]
\centering
\color{revisionblue}
\bitablecaption{네 검토자 원 판정의 재집계(27장면).}{Reanalysis of four reviewers' original labels on 27 scene groups.}
\label{tab:agreement}
\footnotesize
\begin{tabularx}{\columnwidth}{Ycccc}
\toprule
항목 & $\kappa$ & AC1 & 쌍별 일치율 & 불일치 \\
\midrule
시간 요구 & 0.203 & 0.627 & 0.698 & 14/27 \\
행동열 & 0.291 & 0.323 & 0.340 & 23/27 \\
텍스트 일관성 & 0.027 & 0.301 & 0.407 & 23/27 \\
\bottomrule
\end{tabularx}
\end{table}
\begin{table}[t]
\centering
\color{revisionblue}
\bitablecaption{검토자별 텍스트ㆍ시간 요구 판정 분포(각 27건, 익명 A--D).}{Per-reviewer textual and temporal-requirement label distributions, 27 ratings per field.}
\label{tab:reviewer-distribution}
\footnotesize
\begin{tabularx}{\columnwidth}{Ycccc}
\toprule
판정 & A & B & C & D \\
\midrule
텍스트 모순 없음 & 9 & 18 & 10 & 25 \\
텍스트 모순 & 10 & 0 & 1 & 1 \\
외부 근거 필요 & 4 & 4 & 5 & 0 \\
문장 모호 & 1 & 0 & 0 & 0 \\
해당 없음 & 3 & 5 & 11 & 1 \\
\midrule
시간 요구 있음 & 18 & 22 & 15 & 26 \\
시간 요구 없음 & 9 & 5 & 11 & 1 \\
시간 요구 모호 & 0 & 0 & 1 & 0 \\
행동열의 서로 다른 라벨 수 & 20 & 15 & 18 & 10 \\
\bottomrule
\end{tabularx}
\end{table}

표~\ref{tab:agreement}에서 텍스트 일관성의 $\kappa$는 0에 가까웠고 대부분의 장면에서 네 명의 판정이 달랐다. 표~\ref{tab:reviewer-distribution}의 검토자별 분포도 모순이나 해당 없음을 고르는 경향이 크게 달라, 라벨 적용과 해석이 불안정했음을 보여 준다. 그 원인이 전문성이나 교육의 차이라고 확정할 수는 없으며, 분석 단위는 27장면이므로 판정 수를 독립된 정답 수로 세지 않는다.

AC1 추가 재분석은 세 항목 모두에서 Fleiss 계수보다 높았다. 범주 편중에 따른 보정 민감도이며 지침의 안정성이 개선됐다는 뜻이 아니다. 특히 시간 요구는 ``있음''에 판정이 몰려 있어 같은 쌍별 일치율에서도 두 계수의 우연 일치 기대값이 크게 달랐다. 텍스트는 원시 일치 자체가 낮으므로 계수를 바꾸어도 해석이 달라지지 않는다. 기존 $\kappa$를 대체하지 않고 함께 보고한다.

원 메모에는 시간 요구가 없는 행동 조절, 해제 조건과 행동 순서의 구별, 선행 관계를 연결할 정보 부족 같은 해석이 남아 있다(표~\ref{tab:disagreement-examples}). 이는 개별 의견이며 합의 원인을 전수 분석한 것이 아니다. 합의 기록은 모두 모순 없음이지만 사유 기록이 없으므로 이를 객관적인 자연 정답으로 쓰지 않는다.

\begin{table*}[t]
\centering
\color{revisionblue}
\bitablecaption{원 판정의 주요 불일치 예시. 개별 메모는 합의 사유가 아니다.}{Examples of original disagreements; individual notes are not consensus rationales.}
\label{tab:disagreement-examples}
\scriptsize
\begin{tabularx}{\textwidth}{lL{.29\textwidth}Y}
\toprule
익명 장면 & 네 원 판정(A/B/C/D) & 보존된 개별 의견의 요지 \\
\midrule
G01 & 모순/해당 없음/해당 없음/모순 없음 & C: 속도 조절 외 시간 요구가 없다고 해석 \\
G04 & 모순/외부 근거 필요/외부 근거 필요/모순 없음 & C: 신호가 바뀐 뒤 출발했는지 알 수 없어 증거 필요 \\
G22 & 모순 없음/모순 없음/외부 근거 필요/모순 & C: 보행자 통과 여부에 외부 근거 필요; D: 정지ㆍ양보 요구로 해석 \\
\bottomrule
\end{tabularx}
\end{table*}

\subsection{궤적과 보조 반응 결과}
\label{sec:trajectory-results}
문장끼리의 일관성과 별개로 CoC가 요구한 행동이 기록된 자차 움직임과 맞는지도 확인했다. 그러나 대부분의 장면은 판단할 근거가 부족했고 반응 판정은 탐지 설정과 제한 시간에 따라 달라졌다.

27장면의 궤적 비교는 일치 2, 불일치 7, 미상 18이며 보존된 결과를 다시 실행해도 장면별 판정은 같았다. 표~\ref{tab:trajectory-cases}와 부록 표~\ref{tab:trajectory-details}의 사유에서 보듯 미상의 주된 원인은 지원하지 않는 행동과 단계 시점ㆍ해제 근거의 부족이었다. 미상 장면 중 다수에는 일치 사건도 있었지만 다른 사건의 근거가 부족해 장면 판정을 확정하지 못했다. 따라서 현재 자료에서 궤적 대조를 확정할 수 있는 범위는 지원 행동과 관측 행동을 해석할 수 있고 단계 시점이 주어진 사건으로 제한된다. 불일치 장면은 문장 사이의 모순이나 위험 주행의 독립 정답이 아니다.

\begin{table}[t]
\centering
\color{revisionblue}
\bitablecaption{27장면 궤적 비교 요약. 사유별 장면 수는 중복된다.}{Summary of the 27-scene trajectory comparison; reason counts overlap.}
\label{tab:trajectory-cases}
\footnotesize
\begin{tabularx}{\columnwidth}{lcY}
\toprule
장면 판정 & 장면 수 & 포함된 사건의 사유(장면 수) \\
\midrule
일치 & 2 & 모든 사건에서 표현한 행동과 종방향 움직임이 일치 \\
불일치 & 7 & 종방향 행동 불일치(7), 일치 사건 동반(4); 문장 모순이나 안전 위반의 정답은 아님 \\
미상 & 18 & 지원 행동 부재(14), 단계 시점ㆍ해제 근거 부족(9), 구문 해석 미상(6), 관측 행동 불명확(1), 일치 사건 동반(10) \\
\bottomrule
\end{tabularx}
\end{table}


행동 방향이 분명한 보조 반응 분석 대상 134건은 기본 설정에서 반응 확인 125건, 사건 시점에 이미 만족 4건, 제한 시간 내 미확인 5건이었다. 탐지 대상 130건을 27설정에서 재집계하면 항상 검출 72, 설정 의존 57, 항상 미검출 1이다. 확인되지 않은 5건은 제한 시간을 4초 이상으로 늘리면 반응이 검출되는 조합이 크게 늘었으나(5사건$\times$27설정 중 $D=3$초 18개, $D=4$초 80개, $D=5$초 103개), 한 장면은 5초에서도 기본 설정으로 검출되지 않았다. 반복 판단마다 시계를 다시 시작하면 두 장면의 3초 판정이 통과로 바뀌었다. 이 수는 탐지 설정의 집계이며 UPPAAL 질의를 통과한 건수가 아니다. 이 결과는 판정에 사용한 설정을 함께 기록해야 함을 보여 주며 차량 안전 기준을 세우지 않는다.
\end{revision}
